\documentclass{article}
\usepackage{amsmath,amssymb,mathrsfs}
\usepackage[margin=1in]{geometry}
\allowdisplaybreaks
\usepackage[UTF8]{ctex}
\begin{document}

\section*{去中心化 H-合约：具有内生性零售商对冲功能的生产者}

\begin{equation}
\Pi_{\mathrm P}^{\mathrm H}
=\max_{w_\tau,\,\lambda\geqslant0}
\mathbb E_0^{\mathbb Q}\left[
(w_\tau-c_\tau)
\left(\frac{\bar A_\tau-w_\tau(1+\lambda)}{2\xi}\right)^+
\right].
\tag{12}
\end{equation}
\begin{equation}
\text{subject to}\qquad
\mathbb E_0^{\mathbb Q}\left[
 w_\tau\left(\frac{\bar A_\tau-w_\tau(1+\lambda)}{2\xi}\right)^+
\right]\leqslant B.
\tag{13}
\end{equation}
(12)中的决策是$\mathscr F_\tau$可测量的批发价格菜单$w_\tau$以及标量$\lambda$在时间零点的选择。打印出的领域是$w_\tau,\lambda\geqslant0$；命题6的证明也明确使用了$0\leqslant w_\tau$。乘数$\lambda$源自零售商的预期预算限制，而不是生产者的路径约束。在(13)之后的文字表明，在最优解的情况下，如果最优$\lambda$大于零，那么(13)就是紧的；这种情况只会在预算限制具有约束性时才发生。这是一个最优性声明，而不是对(12)-(13)中每个可行选择所施加的额外等式。

\subsection*{固定的输入、信息和经济定义保持不变。}

采购时间$\tau$在此公式中是固定的。生产商处于斯塔克伯格领导地位；零售商则对提供的价格范围做出回应。谈判在$t=0$进行，采购和生产在$\tau$进行，而零售结算则在$T$完成。共同的输入限制包括：
\[
T>0,\qquad 0\leqslant\tau\leqslant T,\qquad A\geqslant0,\qquad \xi>0.
\]
这里，$A$是随机的市场规模参数，$\xi$是一个已知常数，$B$是零售商固定的初始预算，而$c_\tau$是在固定采购时间下的单位生产成本常数。论文假设$c_\tau$随着$\tau$的增加而增长，且生产能力是无限的。所有模型参数都是常识性的。两个代理人都是风险中性的，并且使用相同的已选定的等价鞅测度$\mathbb Q$来评估收益。

在概率空间$(\Omega,\mathscr F,\mathbb Q)$上，过滤$\{\mathscr F_t\}_{0\leqslant t\leqslant T}$是由可交易的安全-价格过程$X_t$生成的。随机变量$A$是$\mathscr F$-可测量的，但不是$\mathscr F_T$-可测量的；特别是，它是$\mathscr F_T\ne\mathscr F$-可测量的。金融回报和$A$是相互依赖的。关于$A$在采购之前的信息，仅有的就是从$X_t$在$[0,\tau]$期间获得的。条件期望值和需求预测是由……定义的。
\[
\mathbb E_s^{\mathbb Q}[\,\cdot\,]
=\mathbb E^{\mathbb Q}[\,\cdot\mid\mathscr F_s],
\qquad
\bar A_\tau:=\mathbb E_\tau^{\mathbb Q}[A].
\]
假设1是：
\[
\mathbb E_\tau^{\mathbb Q}[A]\geqslant c_\tau
\qquad\text{for all }\tau\in[0,T].
\]
附带的文字说明指出，这一条件在任何状态下，在时间$\tau$时都是成立的。

整笔付款是线性的，而清算价格和运营收益则是
\[
\mathscr W(q)=wq,
\qquad\text{clearance price}=A-\xi q_\tau,
\]
\begin{equation}
\Pi_{\mathrm P}:=(w_\tau-c_\tau)q_\tau,
\tag{1}
\end{equation}
\begin{equation}
\Pi_{\mathrm R}:=(A-\xi q_\tau)q_\tau-w_\tau q_\tau.
\tag{2}
\end{equation}
数量$w_\tau$、$q_\tau$以及零售商的经济索赔$G_\tau$都是可以用$\mathscr F_\tau$来衡量的。符号$q_\tau=q(w_\tau)$表示对所提供的菜单的回应；作者明确不要求$q_\tau$能够用由$w_\tau$生成的较小sigma代数来衡量。

\subsection*{财务可接纳性以及内生预算}

存在一个无风险的现金账户，其利率为零。因此，$X_t$是一种$\mathbb Q$-martingale策略。如果选择该文档的连续时间版本，那么$\mathscr F_t$-可预测的自筹资金策略的收益为
\begin{equation}
G_\tau(\theta):=\int_0^\tau\theta_s\,dX_s,
\tag{3}
\end{equation}
其中，$\theta_s$表示在时间$s$时持有的安全单位数量；  
自我融资决定了现金账户的状态。在本文未明确规定的关于可行策略的约束条件下，$G_t(\theta)$是一个$\mathbb Q$-鞅。

完整的金融市场假设认为，每一个可以积分的、可测量的$\mathscr F_\tau$型索赔，都可以通过一种可预测的$\mathscr F_t$型自我融资策略来实现。因此，本文对索赔进行了优化处理，并忽略了$\theta$这一参数。其对收益过程的规定是：
\[
G_s:=\mathbb E_s^{\mathbb Q}[G_\tau]
\quad\text{is a }\mathbb Q\text{-martingale for }s<\tau,
\qquad
\mathbb E_0^{\mathbb Q}[G_\tau]=G_0.
\]
对于选定的 H-合约，初始交易资本为零，因此
\[
G_0=0,\qquad \mathbb E_0^{\mathbb Q}[G_\tau]=0,
\qquad B_\tau:=B+G_\tau,
\]
\[
w_\tau q_\tau\leqslant B+G_\tau=B_\tau
\qquad\text{for all }\omega\in\Omega.
\]
零售商可以选择索赔金额以及采购数量。  
生产者不进行交易。关于技术可集成性和过滤条件，本文未作详细说明，参考了Duffie（2001）的研究；此处未提供外部参考，且未经验证。也没有施加任何额外的可接受性标准。

\subsection*{继承的零售商问题及其解决方案：命题6}

对于固定的菜单$w_\tau$，零售商会选择$\mathscr F_\tau$可测量的$q_\tau$和$B_\tau$：
\begin{equation}
\Pi_{\mathrm R}^{\mathrm H}(w_\tau)
=\max_{q_\tau\geqslant0,\,B_\tau}
\mathbb E_0^{\mathbb Q}
\left[(\bar A_\tau-\xi q_\tau-w_\tau)q_\tau\right]
\tag{9}
\end{equation}
\begin{equation}
\text{subject to}\qquad
w_\tau q_\tau\leqslant B_\tau
\qquad\text{for all }\omega\in\Omega,
\tag{10}
\end{equation}
\begin{equation}
\mathbb E_0^{\mathbb Q}[B_\tau]=B.
\tag{11}
\end{equation}
这些行定义了“追随者问题”，而(12)中的响应就是从这个问题中得到的；它们并不是第二个生产者优化问题。

命题 6 的定义如下：
\[
\mathscr Q_\tau:=\left(\frac{\bar A_\tau-w_\tau}{2\xi}\right)^+,
\qquad
\mathscr X:=\{\omega\in\Omega:B\geqslant\mathscr Q_\tau w_\tau\},
\qquad
\mathscr X^c:=\Omega-\mathscr X.
\]
这两种情况都是有条件的替代方案，而非同时存在的限制条件。

\paragraph{案例 1.}
如果
\[
\mathbb E_0^{\mathbb Q}[\mathscr Q_\tau w_\tau]\leqslant B,
\]
然后
\[
q(w_\tau)=\mathscr Q_\tau.
\]
该命题指出，存在无限多的最优解。其自然的选择是……
\[
G_\tau=[\mathscr Q_\tau w_\tau-B]\,
\begin{cases}
\delta,&\omega\in\mathscr X,\\
1,&\omega\in\mathscr X^c,
\end{cases}
\qquad
\delta:=
\frac{\displaystyle\int_{\mathscr X^c}
 [\mathscr Q_\tau w_\tau-B]\,d\mathbb Q}
{\displaystyle\int_{\mathscr X}
 [B-\mathscr Q_\tau w_\tau]\,d\mathbb Q}.
\]
此声明是一种可选的最优选择，并非要求每个最优对冲策略都必须采用这种形式。所显示的比例没有规定的零分母约定；这里也没有添加任何零值。

\paragraph{案例 2.}
如果
\[
B<\mathbb E_0^{\mathbb Q}[\mathscr Q_\tau w_\tau],
\]
然后
\[
q_\tau(w_\tau)=
\left(\frac{\bar A_\tau-w_\tau(1+\lambda)}{2\xi}\right)^+,
\qquad
G_\tau=q_\tau(w_\tau)w_\tau-B,
\]
其中，$\lambda\geqslant0$负责解决
\[
\mathbb E_0^{\mathbb Q}\left[
 w_\tau\left(\frac{\bar A_\tau-w_\tau(1+\lambda)}{2\xi}\right)^+
\right]=B.
\]
在这种情况下，等式属于零售商的条件响应结构；它并不取代不等式（13）在整个生产者可行集中的存在性。

\paragraph{原文证明不一致。}
命题6的证明中，补充松弛度行被打印出来了。
\[
\beta_\tau(w_\tau q_\tau-B+G_\tau)=0,
\qquad \beta_\tau=\lambda,\quad \beta_\tau\geqslant0.
\]
$G_\tau$中的符号与(10)以及$B_\tau=B+G_\tau$中的符号存在冲突。这一行证明内容被保留下来，只是为了指出印刷上的不一致之处；它并不作为额外的约束条件来约束(12)--(13)。上述关于操作对冲和预算的定义遵循了(9)--(11)以及命题6的陈述。

\subsection*{链接解决方案的特征描述：命题7}

命题7描述了如何求解(12)--(13)式，而不是用另一种生产者模型来替代其可行集。让$\phi^{\mathrm H}$成为能够解决该问题的$\phi\geqslant1$的最小值。
\[
\mathbb E_0^{\mathbb Q}\left[
\left(\frac{\bar A_\tau^2-(\phi c_\tau)^2}{8\xi}\right)^+
\right]\leqslant B,
\qquad
\delta^{\mathrm H}:=\phi^{\mathrm H}c_\tau.
\]
在固定的采购时间内，$\delta^{\mathrm H}$并非随机变化的。最优菜单和订购数量能够满足需求。
\begin{equation}
w_\tau^{\mathrm H}=\frac{\bar A_\tau+\delta^{\mathrm H}}{2},
\qquad
q_\tau^{\mathrm H}=
\left(\frac{\bar A_\tau-\delta^{\mathrm H}}{4\xi}\right)^+.
\tag{14}
\end{equation}
条件性预期收益为：
\begin{equation}
\begin{aligned}
\Pi_{\mathrm P\mid\tau}^{\mathrm H}
&=\frac{(\bar A_\tau+\delta^{\mathrm H}-2c_\tau)
(\bar A_\tau-\delta^{\mathrm H})^+}{8\xi},\\
\Pi_{\mathrm R\mid\tau}^{\mathrm H}
&=\frac{\big((\bar A_\tau-\delta^{\mathrm H})^+\big)^2}{16\xi}.
\end{aligned}
\tag{15}
\end{equation}
收益表示遵循该论文的惯例：
\[
\Pi_{\mathrm P}^{\mathrm H}
=\mathbb E_0^{\mathbb Q}[\Pi_{\mathrm P\mid\tau}^{\mathrm H}],
\qquad
\Pi_{\mathrm R}^{\mathrm H}
=\mathbb E_0^{\mathbb Q}[\Pi_{\mathrm R\mid\tau}^{\mathrm H}].
\]
所述的关闭特性如下：
\[
q_\tau^{\mathrm H}=0
\quad\Longleftrightarrow\quad
\bar A_\tau\leqslant\delta^{\mathrm H}.
\]
附录中命题7的证明（物理PDF第17-18页）表明，$\lambda$的最优值为零。这是一个解决方案的结论，而不是对(12)的额外限制。附录中的停止时间泛化方法和中间松弛优化问题，并未被所选的固定时间公式所替代。

\paragraph{印刷的单调性不一致性。}
在 (15) 之后，论文指出 $\delta^{\mathrm H}$ 在 $B$ 中是“非递减的”。这一点与上述的最小阈值定义相矛盾，同时也与接下来的句子相矛盾，即现金约束更严格的零售商将面临更高的批发价格。这里的印刷文字被记录下来，但它并不是额外的模型条件。

\paragraph{来源范围。}
操作来源是 Caldentey 和 Haugh 的《\emph{供应合同与金融对冲}》，《运筹学》57(1), 47--65 (2009)中的第2节和第4.1节。$\bar A_\tau$的定义见第3.1节。附录中的第7个命题的证明也包含在内。物理PDF文件中的第12--14页以及第17--18页包含了相关的解答和后续内容；这些内容的提供无需任何额外的说明。F-合同、集中式模型、合同比较以及辅助性优化问题都不属于所选的论述范围。

\end{document}
